Future development of the engineering tool could bring additional features and improvements 
to expand its capabilities, more advanced electrical protection analysis of the diagrams. 
Future development of the software should focus on expanding its engineering capabilities, 
improving automation, and enhancing the overall user experience. A potential improvement is 
in the integration of Artificial Intelligence (AI) for the automated diagram recognition. 
This functionality could allow users to import existing PDF engineering schematics and 
automatically identify electrical components, symbols, labels, and connections without the 
additional need to click a verification button. The recognized components could then be 
converted into editable objects within the Sandbox, reducing the time required to manually 
recreate existing diagrams.

Additional Future Improvements:
\begin{itemize}
  \item Automated Topology Validation - Improve the detection of disconnected components, incorrect connections, and incomplete electrical component configurations.
  \item Improved Analysis Calculation Capabilities - Expand beyond momentary DC battery load calculations to include additional protection and electrical system analyses. 
  \item Improved PDF Importing- Enhance background alignment, scaling, symbol recognition, and 
diagram conversion capabilities.
    \item Expanded Component Library - Include additional substation equipment, protective relays, 
transformers, circuit breakers, and electrical devices. 
    \item Automatic Report Generation - Generate engineering reports containing one-line diagrams,
calculation results, protection-zone information, and identified system issues.


\end{itemize}
